Tunnel boring machines (TBMs) are widely used in long and deep tunnel construction,
where excavation responses emerge from the coupled effects of geological conditions,
machine geometry, and operational process. Under fractured, water-rich, high-stress, or
locally deformable ground conditions, the interaction state between surrounding rock
and the TBM may change rapidly and manifest through abnormal thrust, torque,
penetration, advance rate, or shield-related responses. These monitoring variables are
informative external expressions of rock--TBM interaction, but the geological body and
the cutterhead or shield form a spatial relation that evolves as excavation advances.
Monitoring curves alone provide limited evidence about where the interaction occurs,
which machine components are implicated, and how local interaction relations evolve
along chainage.

Most existing methods process TSP-derived geological indicators and TBM parameters as
tabular variables or time-series variables. Such representations can support trend
prediction, but they cannot answer which rock spatial unit forms a candidate interaction
relation with which machine component at a given excavation step
\citep{yagiz2008utilizing,hassanpour2010tbm,armaghani2018prediction,hochreiter1997long}.
Sequence models preserve temporal dependence but treat excavation as an ordered signal
rather than a spatially structured interaction process, so rock location, machine
component structure, and local rock--machine candidate relations are not retained within
a unified representation.

Rather than pursuing a more complex prediction model, this study constructs a
chainage-referenced rock--TBM interaction graph sequence that organises rock voxels,
TBM surface samples, and operational response sequences into a computable
spatio-temporal graph. TSP-derived geological attributes are formalised as rock voxel
nodes, while the cutterhead and shield are discretised as parameterised TBM surface
nodes. Candidate rock--machine edges are screened through spatial proximity,
signed-normal consistency, active-zone filtering, and excavation-state constraints.
Instead of interpreting edge weights as contact-force measurements or relying on
sparse jamming labels, the framework uses continuously recorded monitoring responses to
supervise the relative importance of geometry-screened candidate relations, and
aggregates the learned relevance at edge, surface-node, component, and chainage
levels to generate response-consistent diagnostic views.

The main contributions are threefold. First, the study proposes a chainage-referenced
rock--TBM spatial representation that formalises heterogeneous spatial entities---TSP
voxels, TBM surface components, and chainage-indexed monitoring responses---into a
unified entity model preserving geometric and topological properties. Second, it
proposes a geometry-constrained rock--machine interaction graph sequence in which
candidate edges are restricted by active-zone membership, distance, surface-normal
compatibility, and excavation state, ensuring that learned relations remain
geometrically plausible rather than unconstrained. Third, it uses response-supervised
prediction to examine the response consistency of the learned graph relations and
outputs edge-, component-, and chainage-level relevance diagnostics, with clear
interpretation limits, as spatially organised diagnostics rather than contact-force or
risk estimates.
